\documentclass[11pt]{article}

% ============================================================
% Response to Reviewer -- Round 22
% Companion to: revised_paper_v18.tex
% "Mortgage Lock-In and the Federal Reserve's Quantitative
%  Tightening Shortfall"
% Sixteen-item report across three tiers.
% Point-by-point drafted 26 July 2026; Section 7 added 27 July 2026
% recording what changed in the manuscript between drafting and sending.
% NOT SENT.
% Section references are to revised_paper_v18.tex; every numeric
% value is drawn from a committed run artifact or from the code
% path named inline.
% ============================================================

\usepackage[margin=1in]{geometry}
\usepackage{amsmath}
\usepackage{amssymb}
\usepackage{enumitem}
\usepackage{parskip}
\usepackage[hidelinks]{hyperref}

\newcommand{\mlabel}[1]{\texttt{#1}}
\newcommand{\run}[1]{\texttt{#1}}

\title{\vspace{-2em}Response to Reviewer\\[0.3em]
\large Manuscript: \emph{Mortgage Lock-In and the Federal Reserve's
Quantitative Tightening Shortfall}\\[0.2em]
\normalsize Sixteen-item report: Danish counterfactual, fit statistics,
floor form, and estimator scope\\[0.2em]
\footnotesize Point-by-point drafted 26 July 2026; \S7 records what changed in the
manuscript between drafting and sending}
\author{}
\date{27 July 2026}

\begin{document}
\sloppy
\maketitle
\vspace{-2.5em}

Dear Editor and Reviewer,

Thank you for the report. It is the most useful I have had, and I want to be
straightforward about two things before the point-by-point.

First, several items were written against a draft that a prior revision had
already changed. I say so not to deflect but because the reader deserves to know
which of my responses are new work and which are pointers to work already in the
manuscript; I have labelled every item accordingly, and where an item is already
met I quote the sentence that meets it so you can judge for yourself rather than
take my word.

Second, and more important: the items that did land were not marginal. Your
functional-form point (item 3) turned out to have a second half neither of us
had seen, your fit-statistic point (item 2) was right about a defect I had been
reporting for three rounds, and your instinct about the Danish leg (item 1) was
correct on a channel different from the one you named. Chasing them produced
five new pre-committed runs and, along the way, four defects in the manuscript
that no reviewer had raised and my own verification suite could not see. Those
are itemised in \S5, because a report that surfaces errors it did not itself
identify has done its job better than a report that only lists its own hits.

Numbers below are quoted from committed artifacts. New runs follow the project's
spec-before-run convention: a specification commit precedes each result commit,
and each carries parity gates that replay committed values bit-exactly before
anything new is reported.

\section{Item 1 --- the Danish counterfactual}

\textbf{Status: the constructive fix was already in the manuscript; one prong
you raised was live, and it was worse than you argued.}

\subsection{1(a) The floor violation and the U.S.-intercept transplant}

This is implemented, and by exactly the construction you prescribe. The
production leg (\run{danish\_us\_intercept}) keeps the identified flatness but
anchors the level at the U.S. hazard evaluated at zero rate gap, with the
involuntary floor retained: \texttt{competing\_risks.py} passes a zeroed rate gap
into \texttt{prepay\_hazard}, so the floor is applied on the Danish leg exactly as
on the U.S.\ leg. It lands at 5.61\% mean Danish CPR --- your own predicted
``roughly 6\% CPR'' --- and turns the gap positive: $+\$61.2$ billion at the
in-sample calibration point and $+\$28.2$ billion at the off-window calibration
the headline marginal uses.

The Danish-level leg survives only as a labelled bracketing case, and your
diagnosis of it is correct: it prepays at 3.39\%, below the 4\% floor both
frameworks impose. That inconsistency is now disclosed wherever the
$-\$99.9$ billion figure appears, including the Introduction and the Conclusion,
which previously carried the number without the caveat.

\subsection{1(b) ``No Danish mechanism'': the numbers are right, the inference
is not}

Your arithmetic checks exactly --- I verified it from the two parquets rather
than from the tables. The rule-only Danish leg and the $\beta_1 = 0$ null agree
to $0.000107$ percentage points of mean CPR and to $\$0.004$ billion of trapped
balance.

But that identity is forced, and it is forced by the fix you ask for in 1(a). The
gap enters the hazard only through $(-\beta_1) \times \text{gap}$; setting the gap
to zero and setting $\beta_1$ to zero are the same operation on that term. A
gap-flat hazard at the U.S.\ zero-gap intercept \emph{is} the null leg with the
floor retained. So the coincidence is not evidence that the mechanism is absent;
it is what a correct transplant must produce. The manuscript states this
directly: the two are ``one object measured on two accounting legs.''

\subsection{1(c) The discount-rate table --- you were right, and about a larger
channel than you named}

You read the exactly-zero deltas as evidence that market-value repurchase is
absent rather than as immunity to a pricing critique. That reading is correct,
and pursuing it found two things.

The smaller one is that the diagnostic had executed on the \emph{wrong leg}.
\texttt{danish\_discount\_bound.py} contains no call to
\texttt{set\_danish\_moving\_anchor}, so it ran at the module default
(\texttt{dk\_level}); its artifact records a mean Danish CPR of 3.39\%, the
bracketing anchor's speed, not the production leg's 5.61\%. The table caption
nevertheless described it as the production leg and extended the result to the
production gap. Caption and text now name the leg that ran, and state that the
production anchor's PV-independence is structural rather than run-verified.

The larger one is the accounting. The Danish leg credits every retired balance at
\emph{par}: the dynamic-balance loop books roll-off as balance times drain with
no price term anywhere in it. Netting the roughly \$199 billion scheduled
component from the leg's \$669.3 billion of Danish roll-off leaves on the order
of \$470 billion credited at par, against a proxy discount of 32--34\% at the
representative 3.0\%/6.8\% state. A market-value credit would reduce Danish cash
receipts by more than the $+\$61.2$ billion gap itself, and in the direction that
shrinks it. I have not restated the gap under such a credit, because the pipeline
does not distinguish a cash haircut from a balance adjustment; I record it as the
largest un-priced channel in the counterfactual and retract the claim that the
pricing critique attaches to the superseded variant alone.

\subsection{1(d)--1(e) The sweep, and the Berger citation}

The sign-fragility framing you object to is gone: the figure and the text now say
that the anchor, not the refinance channel, decides the sign. Two further
corrections came out of re-reading that figure, both reported in \S5.

On your falsifier --- whether Berger et al.\ estimate 3.2\% per year as a
general-equilibrium prediction for U.S.\ households or as a Danish descriptive
level --- I checked the source directly. It is the latter: the January 2026 draft
uses 3.2\% as ``the unconditional average moving rate among Danish FRM
borrowers,'' a calibration target in its \S4.6. The manuscript's slope-versus-level
distinction is therefore the correct reading of that paper, and it now carries
locators so you can check it in one step.

\section{Item 2 --- recovery percentage as the primary fit statistic}

\textbf{Status: implemented, with one correction to the argument.}

You are right that recovery is an affine transform of cumulative error with a
large additive constant, and that this flatters every estimator. \mlabel{tab:bases}
now carries a terminal cumulative runoff error column, the note states the
identity explicitly, and errors are given as shares of the \$652.8 billion of
runoff actually realized. Your arithmetic reproduces exactly.

One correction. Nothing changes sign under the transform: recovery
$= 1 + \text{error}/764.7$ is monotone. What flips is the \emph{referent} --- an
over-prediction of trapped liquidity is an under-prediction of runoff --- and the
denominator. I have written it that way because the sharper claim is available
without the stronger one: Path B's in-sample $+\$53.8$ billion is simultaneously
an 8.2\% under-prediction of realized runoff, and the headline off-window leg,
which reads as a worse fit at 91.3\%, in fact carries the \emph{smallest} error in
the table on the standalone scorer, $+\$2.8$ billion or 0.4\%.

Your asymmetry point about WSHOMCB rescaling is a correct reading of the code and
now underwrites a different result; see item 6.

\section{Item 3 --- the floor's functional form}

\textbf{Status: the level-versus-form test was already in the manuscript; the
prong you identified as still open was open, and it refuted a claim I was making
implicitly.}

The additive competing-risks form had been run at the production floor and at the
three off-window anchors, giving a form-conditional identified range that stood at
the time at $+3.9$ to $+13.1$ points. That range has since widened, against me, and
I take it up in \S7. Your specific charge --- that the concave-transform check had
been executed inside the censoring regime it needed to escape --- was correct:
\run{concave\_marginal} ran under the hard-maximum form only.

Running the joint cell (\run{concave\_additive\_marginal}) leaves the range intact:
the concave $\times$ additive marginal is $+9.25$ points at the production floor and
$+9.22$ at the headline anchor, inside the hull at both. But it refutes something
the manuscript had been asserting without measurement. \mlabel{tab:uncertainty}
listed the form adjustment and the transform adjustment side by side, which
presumes they add. They do not: the interaction is $-1.47$ points at the headline
anchor, past the $\pm 1$-point convention used throughout. The mechanism is the
censoring at issue --- the concave transform costs $-0.52$ points under the max
form, which absorbs most of it into the floor, against $-1.99$ under the additive
form, which never censors. The table now reports the joint cell and the
interaction rather than implying additivity.

I have also disclosed what the additive form costs in level terms, which the
manuscript had priced only at the production floor: at the off-window anchors the
additive central leg recovers 55.9\% of the benchmark against a 44.7\% null, and
the hull's upper endpoint comes from a cell recovering 61.2\%. The
form-conditional range is therefore a range over specifications of sharply
different aggregate fit, and its upper end is not an equally credentialed member.
That cuts against the additive form rather than for it, and it belongs in the
open.

\section{Item 4 --- the ABM; item 5 --- Path A; item 16 --- scale}

\textbf{Item 4: partly already met, one prong newly closed.} The abstract already
reported the fifty-seed mean rather than the seed-42 draw. Your remaining charge
--- that the cross-design threshold test was a single draw from an estimator with
material seed dispersion --- was correct, and \run{cross\_design\_seeds} now runs
fifty seeds across both randomization layers. The verdict survives and
strengthens: mean 60.2\%, standard deviation 3.3 points, 95\% band 55.1--66.0\%,
the committed draw at the 48th percentile of its own distribution, and all fifty
seeds classifying as undercutting. The band is reported wherever the point
appears.

That sweep also separates the two variants on the manuscript's own admissibility
test, in the direction that was not convenient for me: the recalibrated variant's
deep-gap floor is inside the observed 4--5\% band on all fifty seeds, and the
frozen variant's is outside it on all fifty. Applied consistently, my own test
disqualifies the leg I had been reporting as corroborating.

\textbf{Item 5: your premise does not hold, and I want to be exact about where.}
The abstract's 97.9\% is Path~B's in-sample recovery, not Path~A's; the abstract
contains no Path~A figure. But your underlying concern was sound at one site:
Section~I carried Path~A's 112.4\% describing it as corroboration, and a
bias-respecting sign test (\run{patha\_sign\_test}) does not reject
$\beta_g \leq 0$ under any construction. The corroboration language is retracted
there and in the interpretation section, and the diagnostic table now labels each
row's specification --- it had been interleaving spec v4 and spec v3 rows without
saying so.

\textbf{Item 16: correct, and unmet until now.} The scale check exercised the
household mechanic in isolation, without the DTI wall, the freeze trait,
curtailment, multi-vintage weighting or the settlement kernel, and ran 3.9 points
hot against the production ABM. The claim is now scoped at all five sites to what
was measured. A production-specification rerun is reported separately; its first
execution halted at a pre-committed parity tier on upstream data revision, and I
have recorded the halt rather than relaxing the tolerance.

\section{Defects this report surfaced indirectly}

These were found while chasing the items above. None was raised in the report and
none was visible to the verification suite, which pins numeric literals and
therefore cannot see a false statement about a magnitude, an omitted summary
statistic, or a caption describing the wrong object.

\begin{enumerate}[leftmargin=1.4em]
\item \textbf{A false claim in the sweep figure's caption.} It asserted that both
anchors stay ``an order of magnitude below the superseded \$925.5 billion
extrapolation across the entire sweep.'' The body text five lines earlier gives
$+\$1{,}038.9$ billion at the ceiling, and both artifacts agree; both anchors
\emph{exceed} \$925.5 billion there. Corrected, with both ceilings quoted.

\item \textbf{Negative trapped balances in the same sweep.} At refinance
contributions of 12\% and above the Danish leg's trapped balance goes to
$-\$14.4$, $-\$158.0$ and $-\$290.0$ billion. That is not an economically
meaningful quantity, and the upper half of the axis is now marked directional only.

\item \textbf{Selective reporting in the Path A bootstrap.} The manuscript
reported the median, the interquartile range, the interval and the fraction of
draws above the benchmark, and omitted the mean: \$598.9 billion, or 78.3\% of
the benchmark --- below the paper's own mechanical null. Now reported, with the
reason.

\item \textbf{An ablation that had never been run.} The burnout ablation was
quoted at four sites with no artifact, script or gate behind it, and its figures
reproduce the \emph{minimum} over 100 replicates of an unrelated origination-time
permutation experiment. Executed properly (\run{burnout\_ablation}), the true
value is $-\$6.58$ billion against the claimed $-\$6.6$ billion: the number was
right by coincidence. The run also showed the ablation halves at the headline
calibration, where three interpretive sites had been quoting the in-sample figure.
\end{enumerate}

\section{Items 6--15, and what I have not done}

Item~6's global validation statement has been moved into Section~I. The
loan-level bootstrap you ask for existed, but you were right that it was the
wrong scheme: it resampled within strata and held the 130 strata fixed, so
between-stratum variance was never drawn. Drawing the strata
(\run{bootstrap\_pathb\_cluster}) widens the marginal's interval by roughly a
factor of 37, to $[+4.63, +6.92]$ points. The floor-read layer remains the wider
of the two, so the headline is unaffected --- but ``the single stratified draw
contributes essentially no uncertainty'' was a statement about the scheme, not
the estimator, and it is retracted.

Item~14's third pre-commitment gap was the sharpest thing in the report. The
admissibility precondition excluding the externally anchored variant does not
appear in that run's specification commit, which instead states that the variant
``is evaluated against the SAME pre-registered bands'' and pre-registers the floor
diagnostic as reported ``whichever direction it moves''; the frozen artifact
classifies the result as undercutting. The manuscript now states the
pre-committed classification first and labels the qualification as post hoc, and
the verification suite reads the artifact's classification so a green suite can no
longer pass over it.

Items~7, 8, 10, 11, 12(b), 12(c) and 15 were met in the manuscript already, and I
quote the operative sentences in the point-by-point appendix rather than here.
Item~13's reduction is correct and is now stated explicitly: the no-trade-off
claim reduces to the marginal being bounded by the excess of the null hazard over
the floor, and to zero where the floor binds --- a property of the production form
that lifts under the additive one.

One item I have \emph{not} implemented, with a reason. Item~12(a)'s Ginnie overlay
is promoted into Table~1 but not made the headline level, because \$27.8 billion of
the raw adjustment is a component common to any observed-series overlay; I have
labelled it an accounting-correction variant instead, and I am glad to go further if
you think that understates it.

Item~9 I declined when I first drafted this section, on the ground that cutting the
abstract further meant deleting a disclosed result and that each candidate was a
disclosure a prior round had added deliberately. It stood at 424 words. That
argument was half right, and I have since acted on the half that was wrong: a hedge
only has to travel with the claim it scopes, so when five claims left the abstract
their hedges could leave with them, provided they were pinned somewhere the reader
still meets them. They are --- against the body, by a verification rule written for
the purpose --- and nothing was de-hedged. The abstract went to 263 words. Four
deliberate additions pushed it to 367 words: returning the agent-based model to
it, which the cut had orphaned against a section the model still occupies,
reporting the lock-in contribution as a range rather than as a point (\S7),
stating inside the abstract itself that the range is form-conditional (the
$+3.5$ to $+13.1$ hull travels with it) and what its sampling layer is --- the
wild-cluster-corrected $+2.9$ to $+8.7$ interval on 31 clusters, with the
narrower percentile read named and marked as under-covering --- and marking the
mechanical null's floor as itself read from realized, partly behavioral
turnover, so the null's ``however the rate-responsive margin behaves'' scope is stated honestly. A
subsequent author-directed cut, one addition after it, and round 32's three abstract qualifications
(the mechanical majority labelled as form-conditional with the additive form's null beside
it; the binding layer named as the widest that does have a coverage property, with the
unestimated seasoning ramp's span in the same clause; and the self-refuting ``identifies
levels only'' replaced) had carried it to 427 words; the v20 panel compression, with the re-review round's two hedge additions (the projection-relative clause's upper-bound status; the cost sentence's zero-refinance-edge condition), and the later removal of the abstract's censoring flag once the grid extension below measured the endpoint, has taken it to 263 words: the agent-based-model
and cross-design readings, and then the expectations-benchmark reading, returned to
the body with their hedges, under the same rule --- a hedge travels with the claim
it scopes, and each is pinned against the body by the same verification
machinery --- while every disclosure the panel asked into the abstract (the
wild-cluster inferential status, the form-conditional hull, the behavioral-floor
caveat) survives there in compressed form. The addition is the form condition on the closing policy sentence: the institutional cash-flow cost is small under the production floor form, and the additive form roughly doubles the margin. I would
rather give you the trail than the endpoint, since the endpoint moved in both directions
and the direction of travel is the honest part.

\section{Changes since this response was drafted}

A post-re-review adjudication supersedes one claim below: the few-cluster coverage run's frozen selection rule, applied as written, lands the binding layer on the restricted wild-cluster inversion (coverage 94.3\%/95.1\%) rather than the Webb wild-$t$ (90.8\%/91.4\%, below the pre-committed bar), so the quoted binding interval became $+2.3$ to $+9.1$ points everywhere, with the Webb read (lower edge $+2.9$, upper $+8.7$) retained beside it as a reported rung; the manuscript's ladder note records the adjudication and gate \#125 ties the quoted interval to the artifact's decision field.

Two later changes supersede figures above. The binding interval's lower endpoint was censored at the committed floor-to-marginal grid's 6.0\% end: every rung whose floor-unit lower edge ran past that grid printed the edge row's $+2.3$. A pre-committed grid extension (run \texttt{floor\_grid\_extension}, August 2026: paired legs at 6.25--7.00\%, the 6.0\% row reproduced bit-exactly, the interpolation checked against an independent engine read at 6.91\% to within \$0.005 billion) measures the endpoint instead. The quoted binding interval is now $+1.9$ to $+9.1$ points everywhere; the Webb rung re-prints as $+2.8$ to $+8.7$ under the same extended mapping; the censoring disclosures come out of the manuscript because there is nothing left to disclose; and the lower endpoint now sits below the calibration box's $+2.1$ edge, which the manuscript states without re-anchoring the box. Second, the convolved sampling line above was computed on the demoted Webb rung; re-derived on the adjudicated restricted inversion over the extended grid (run \texttt{layer\_convolution\_restricted}), it prints $[+1.71, +9.38]$, wider than the binding interval as a convolution must be, and the one-number recommendation for sampling error returns to it.


The point-by-point above was written on 26 July 2026 and this letter was not sent
that day. Rather than leave it describing a manuscript that no longer exists, I
record here what has changed. Four of the eight items below move \emph{against} the
paper, and those are the ones I would read first.

\subsection{The Danish buyback credit is now priced, and it can reverse the gap}

The market-value buyback credit I had recorded as the counterfactual's largest
un-priced channel is now priced as a pre-committed incidence bracket (runs
\texttt{buyback\_credit\_bracket} and \texttt{buyback\_discount\_rederived}).
Under a balance-adjustment reading the $+\$61.2$ billion institutional gap
stands; under a cash-haircut reading it reverses, to $-\$51.0$ billion at a
discount re-derived from the leg's own prepayment path on \$470.7 billion of
early face. That bracket's own discount grid was asserted, never derived; the
re-derivation replaces it with a face-weighted 23.8\%, and the reversal
survives, smaller. The conclusion now states
leg (3)'s sign as incidence-conditional rather than as a modest improvement,
and the trilemma paragraph's dissolution claim is scoped to face accounting.
A resolution note travels with the bracket: the benchmark itself is
face-denominated (SOMA current face against par caps), so the
balance-adjustment reading is the benchmark-consistent one, and the cash
reversal prices the foregone par windfall from moving households --- a
transfer, not a resource cost.

\subsection{The form-conditional hull widened at the bottom}

A pre-committed run (\run{concave\_hull\_band\_ends}) closed the sixteen concave
off-window cells the grid was missing, completing the form $\times$ transform product
at thirty-six of thirty-six. Seventeen of the eighteen concave cells fall inside the
hull I quoted in \S3. The eighteenth does not: concave under the hard-maximum form at
the highest defensible floor and the elasticity band's low edge returns $+3.53$
points, below the $+3.89$ lower edge the log-linear cells had set. The
form-conditional range is therefore $+3.5$ to $+13.1$ points rather than the range
\S3 quotes, and the manuscript carries the wider figure throughout. The upper edge is unchanged,
because the concave transform is signed downward and its highest cell reaches only
$+10.83$.

\subsection{There is no ABM counterpart to the $\beta_1 = 0$ null, and the reason is
structural}

You did not raise this, and you would have. The paper insists that levels do not
identify and that only the central-minus-null differential does, and then conducts the
entire ABM comparison on levels. I tried to build the missing null
(\run{abm\_null\_feasibility}) and it is \textbf{not feasible}, for a reason worth
stating rather than burying: Path~B's elasticity is a separable term over a floor that
is measured independently and survives its removal, whereas the ABM's 4--5\%
turnover floor is not a term in its move rule at all but a \emph{product} of it,
obtained by searching the mobility-desire scale until the penalty-bearing rule returns
4--5\% CPR at an 8\% rate. Zeroing the penalty removes the floor along with the
channel. Holding the scale fixed gives a leg scoring $-259$\% of the benchmark;
re-deriving it gives $+116$\%. The implied marginals are $+270.4$ and $-105.3$ points
--- they disagree in \emph{sign}, and neither is interpretable. I therefore report no
ABM differential and say so where the contrast is drawn. The estimator comparison
remains a comparison of levels, now with its limitation named rather than implied.

\subsection{Every correction to the headline points the same way, and the paper now
says so in one place}

This is the change I expect you to care about most, and it was not on your list. Each
downward correction to the lock-in marginal was disclosed in the section that produced
it and nowhere together: the in-sample calibration returns $+9.2$ points, re-anchoring
the floor off window returns $+5.6$, the Ginnie overlay at that same floor returns
$+4.4$, an age-standardised floor read implies near $+3.8$, and an independent Fannie
Mae read of the same off-window cell (5.52\%) sits above the clean band and brackets
the marginal below $+4.3$. Composing the overlay with the age-standardised floor ---
a composition neither run establishes --- would put it near $+3.0$.

A reader assembling those in their head reaches a conclusion the sections never state,
so the manuscript now performs the assembly itself, in one paragraph, and commits to a
posture: holding the production form and the imported elasticity fixed, every
correction to the floor or to the accounting basis moves the headline down; the
binding uncertainty layer is the floor reads' own sampling error, now quoted after a
wild-cluster small-sample correction at $+2.9$ to $+8.7$ points (run
\texttt{floor\_inference\_correction}: the percentile read, $+3.0$ to $+8.0$ points,
under-covers at 31 clusters, mostly at the top, with CR2/CR3 $t$-intervals agreeing);
the headline sits just below the corrected interval's midpoint while every floor or
basis correction falls in its lower half. A second panel round then found the
assembly's completeness claim false --- three disclosed non-floor variants run
downward (the concave transform, the sub-band elasticity, and the moving-share
bracket, since re-run at the headline floor: $+3.36$ points at $s=0.5$) --- so the
manuscript now lines both directions up and attaches no posture to where $+5.6$
sits: the range, not any interior member, carries the identified content. A Webb-weight re-read of the wild-cluster construction (effective cluster count 5.9) re-prints the interval's lower edge at $+2.9$ --- both endpoints within a tenth of a point of the prior read --- and the small-sample-df secondary reads widen to $+2.3$/$+9.6$ at the widest, agreeing in location. The joint overlay-by-age-standardized
cell has since been run under a pre-committed landing rule: it composes to $+2.9$
points, within $0.1$ of the proportional projection, and the paper now reports it
as the assembly's measured lower member --- where the first joint cell this paper
ran, the concave transform under the additive form, returned an interaction of
$-1.47$ points.

The assembly has since gained a calibrated member. The manuscript quotes Aladangady
et al.'s finding that rate-gap lock-in explains 44\% of the 2021--22 mobility decline
and had never used it; scaling the seasoning baseline in both legs until the non-rate
channel traps $0.56/0.44$ times the rate channel's trapped liquidity puts the root at
75.4 PSA and the marginal at $+0.9$ points at the headline floor (run
\texttt{scaled\_null\_housing\_activity}). The variant documents an upward bias under both forms; under the
production form the corrected member lies below the headline, and above it
under the additive form.

The empirical benchmark's four exact-zero months have been traced to their
mechanism: clip-induced month-boundary allocation artifacts of the SOMA reporting
series (the reported month's rolloff falls below the scheduled-amortization floor;
the mass posts into the next month). The levels-only concession is now stated at
the timing-descriptor sites themselves, not only in the appendix. No timing
exhibit was re-scored and no committed number moves.

The two sampling layers --- the corrected floor read and the loan/stratum
resample --- are now also convolved as independent draws and reported as a
labeled second line, $[+2.80, +8.99]$ points: wider than either input, offered
as the single number for a reader who wants one, and not substituted for the
floor read because the independence is an assumption, not a measurement (the
comonotone bound is disclosed alongside).

The ridge-device verification battery, previously recorded only under spec v3
with a 'not recomputed at adoption' disclaimer, has been recomputed under the
adopted seasonal specification: all three facts survive (grid identity at
machine precision, drop-fit within a pre-committed 0.005, exact reference
invariance of the unpenalized refit), and the disclaimer is deleted at the two
sites this rerun covers. The resampling-distribution disclaimer stays until the
joint bootstrap rerun lands.

A 150-start dispersed-optimization battery on the production panel finds no
penalized objective better than the adopted point (the anchor is evaluated at
the frozen coefficient vector, not refit), while only 1 of 50 primary-leg
starts converges into the production basin --- the manuscript's two
branch-instability sentences now carry both numbers, and the exclusion of
Path A from headline figures is strengthened.

The coupon-convention audit found the full-book reweight matching borrower note
rates against security pass-through coupons. On a single convention the composed
cell reads 98.1\% rather than the committed mixed-basis 100.0\% (both are now
labeled; the mixed-basis figure is retained as the measurement it is), the
coupon-composition bound tightens from two points to $+0.2$, and the identified
marginal is bit-invariant to the entire correction --- structurally, since the
conversion never touches a hazard input.

The back-out's amortization convention is likewise corrected: at the pool's
note-rate WAC the empirical CPR level rises by a deterministic +0.27 points with
every peak lag and U2 unchanged and every trapped-dollar figure bit-identical;
the Theil table is refreshed on the corrected basis and the printed empirical
mean carries its convention label.

The Path A refit-and-resimulate distribution has been rerun with each
replication's full coefficient vector at both specifications. The joint interval
is ten times narrower than the committed one and is not adopted: more than half
its draws --- every rogue-mode replicate among them --- sit at exactly the
zero-voluntary-prepayment floor, measured by a probe through the same
coefficient handoff, so its upper half is mechanical. The committed interval's
-\$4,460B tail and the joint pile are one optimizer pathology rendered two ways,
and the manuscript now says so where the interval is quoted. The point-head
parity leg reproduces the committed points to \$0.0000B at both specifications,
and the spec-v4 leg fails fewer replications than spec v3 (1 vs 2 of 200).

The vintage counterpart of the Ginnie overlay now exists: scoring the
out-of-window vintage share by its observed Fannie speeds in both legs confines
the marginal to the in-window-vintage share at +3.7 points, pure share scaling
(0.66x, placebo within 0.0002). Both overlays are now labelled as what they
are --- changes of estimand, re-scoping the marginal onto a sub-book, not
corrections to the headline.

The within-stratum control three reviewers asked for is not literally
computable (the gradient axis and the stratum id share the coupon dimension),
but its substance is: standardizing the deep bucket to the shallow bucket's
vintage-FICO-LTV composition leaves the gradient at +3.44 points against the
model's implied +0.94, and an exact decomposition puts only +0.77 of the raw
+4.20 on composition. The composition-confounding defence is replaced by the
tested statement; the exhibit stays declined on the narrower surviving ground
(within-window refinancing, and a level gradient is not a dollar marginal),
and the assembly gains its first upward realized-data entry, labelled
directional.

Three things keep that from being a retraction, and all three are in the same
paragraph. The interval does not contain zero, and the sign is fixed by the window's
rate configuration rather than by any of these calibrations. The corrections that are
not floor corrections run the other way: the additive form returns $+11.2$ points at
the same anchors, and the Fonseca--Liu (2024) anchor would put the marginal at $+11.5$
at the production floor. And the three readings I decline elsewhere all point to a
larger channel, not a smaller one.

The front of the paper has been rewritten to match. The abstract, the introduction,
Table~1 and the conclusion now lead with the range and name $+5.6$ inside it, where
they previously led with the point and carried the range afterwards. Two defects
surfaced in the course of that: the introduction had been quoting only the depth-cut
band, $+4.3$ to $+6.8$, as its uncertainty --- the narrower, non-binding layer --- and
Table~1 carried no binding layer at all.

\subsection{The censoring share, and what does not follow from it}

At the off-window floor the paper headlines, the central leg's hazard is pinned to the
turnover floor in 68.8\% of evaluated loan-months, against 36.3\% at the in-sample
calibration point. That is now disclosed where the identification is argued, together
with what does \emph{not} follow from it, because the obvious reading is wrong twice
over: a pinned month's contribution is \emph{truncated} at the floor rather than
erased --- the $\beta_1 = 0$ leg is pinned in only 35.8\% of loan-months at the same
floor, so the two legs still separate wherever the null clears it --- and these are
unweighted loan-month counts rather than balance-weighted shares, so nothing follows
about what fraction of the dollar marginal the uncensored months carry.

\subsection{Two estimators demoted in presentation, with no result touched}

Your items 4 and 5 both pointed at exhibits carrying more prominence than their
evidential standing. I have gone further than the report asked in both cases, and
neither change moves a number.

Section~IV now \emph{leads} with the cross-design result rather than the production
13.6\%: the same decision rule on real Freddie structural covariates recovers 60.2\%
averaged over fifty seeds, above the 50\% threshold fixed before that run at which the
section's own paradigm reading is undercut. The production estimate is named as the
synthetic-population special case, which is what it is. Three bounds travel with that
lead wherever it appears --- the same reweight that lifts the recalibrated leg to
76.3\% at book composition moves the frozen leg \emph{down} to 12.6\%; the covariates
swapped in are structural only, so no leg replaces a synthetic behavioural population
with a real one; and the verdict is that the contrast cannot, on that test alone, be
cleanly attributed to modelling paradigm rather than to calibration and data source.

Path~A occupied three of the paper's first seven tables while its refit-and-resimulate
interval spans zero and its sign claim was retracted in response to your item~5. Its
estimation-panel and bootstrap tables have moved to the appendix verbatim; one
diagnostic table remains in the body. No cell, note or caveat was edited --- in
particular the \$598.9 billion bootstrap mean of \S5 travels with its table.

\subsection{Two further pre-committed runs, both null results}

The benchmark is convention-dependent, and I now measure by how much. Rebasing the
\$764.7 billion cap-shortfall from calendar months to settlement months
(\run{settlement\_months\_benchmark}) moves it to \$722.7 billion, $-5.5$\%, and the
entire shift is mechanical: the QT window is half-open, so cap convolved past the
window's end drops out, and that dropped quantity equals the shift exactly. No dollar
figure in the paper moves --- the marginal is \$42.6 billion on either denominator ---
and the calendar convention remains production, as pre-committed.

The production 43\% DTI threshold sat at the minimum of its own sweep, which reads
badly (\run{dti\_nonmonotonicity\_decomposition}). It is an artifact of the
floor-retention rule, which re-calibrates the mobility scale only where the fixed-scale
floor leaves the observed band, and across those thresholds fires exactly once. Hold
the scale still and the sweep is monotone. No headline moves under either branch,
which was pre-committed.

\subsection{Presentation changes: a designated online appendix and new exhibits}

Four further changes are presentational and move no figure this letter quotes. The
manuscript now designates everything after the conclusion as an online appendix, with
the run ledger, the specification blocks, the notation glossary, and four robustness
units moved there verbatim and their callers re-pointed. Two elasticity-discipline
exhibits landed: the marginal as a curve in the imported elasticity, read down to
half the adopted central value, and a bracket applying the elasticity to an assumed
moving share of the total hazard. And an audit table now censuses every
non-mechanical post-run adjudication, in both directions, so the adjudication record
the manuscript relies on can be checked in one place.

\subsection{Verification}

The suite has grown by roughly a tenth in checks and a third in tests since the
point-by-point above; the current counts are stated in the manuscript rather than
here, because a count quoted in a letter is a number with nothing checking it.

Several of the additions pin \emph{prose} rather than numbers, which is new. The
assembly paragraph and the abstract's ordering are checked as contiguous spans within
a single paragraph, because every literal in them already occurs elsewhere in the
manuscript and a whole-file count would be satisfied by the very scattering they exist
to end; and two of the checks assert an \emph{order} rather than a presence, because
every span survives the arrangement each was written to prevent.

Each new check was mutation-tested, and the exercise earned its keep twice. One
battery found that the check it was written for had failed to pin its own headline
number --- the censoring share, which was the whole point of that check. Another
reported a failure that turned out to be the test's fault rather than the check's: it
had deleted an earlier copy of the span from a different section and left the target
untouched. This letter is now covered too, which is what should have been true when I
first drafted it.

\vspace{1em}
\noindent Thank you again. The report cost me more runs than any previous one and
the manuscript is materially more honest for it.

\end{document}
